\clearpage\SourceLeaf{12}
\begin{center}\bfseries [NOTES ON CUBIC AND BIQUADRATIC RESIDUES.]\end{center}
\begin{center}[1.]\par\smallskip A most elegant observation made by induction.\end{center}
\noindent\hspace*{\parindent}2 is a cubic residue or nonresidue of the prime $p$ of the form $3n+1$, according as $p$ can be represented by the form
\[xx+27yy\qquad\text{or}\qquad4xx+2xy+7yy.\]
\par\noindent\hspace*{\parindent}3 is a residue or nonresidue, according as $p$ can be represented by
\[\begin{aligned} &xx+243yy\quad\text{or}\quad4xx+2xy+61yy\\
\text{or}\quad&7xx+6xy+36yy\quad\text{or}\quad9xx+6xy+28yy.\end{aligned}\]
\par\noindent\hspace*{\parindent}5 is a $\left.\begin{array}{l}\text{Residue}\\\text{Nonresidue}\end{array}\right\}$ if $p$ is represented
\begingroup\small
\[\text{by }\left\{\begin{array}{l}(1,0,675),\ (25,0,27),\ (13,1,52),\ (4,1,169)\\
(7,2,97),\ (9,3,76),\ (19,3,36),\ (25,5,28),\ (25,10,31),\ (27,9,28).\end{array}\right.\]
\endgroup
\begin{center}[2.]\end{center}
\noindent\hspace*{\parindent}A criterion for determining whether the number $m$ is a cubic residue of the prime $p$ of the form $3n+1$.

Let
\[4p=aa+27bb\]
